\documentclass[11pt,a4paper]{article}
\usepackage{fontspec,amsmath,amssymb,amsthm,microtype}
\setmainfont{lmroman10-regular.otf}[BoldFont=lmroman10-bold.otf,ItalicFont=lmroman10-italic.otf,BoldItalicFont=lmroman10-bolditalic.otf]
\usepackage[margin=25mm]{geometry}
\usepackage[unicode,hidelinks]{hyperref}
\hypersetup{pdfauthor={LangevinLean Project}}
\setlength{\emergencystretch}{3em}
\setlength{\parskip}{0.4em}
\setlength{\parindent}{0pt}
\allowdisplaybreaks[2]
\raggedbottom
\title{\texorpdfstring{Cubic velocity and quartic position errors of the Hessian correction}{Cubic velocity and quartic position errors of the Hessian correction}}
\author{LangevinLean Project}
\date{}
\begin{document}\maketitle
Draft mathematical exposition; these notes may contain mistakes.


The free-OU path \(U\) contains the exact initial velocity and the shared Brownian convolution. The exact velocity is \(U\) minus the damped force integral. Correcting the frozen-force approximation by the Hessian acting on the integrated free path leaves two explicit errors: the quadratic force Taylor remainder and a triple force-memory integral. Their orders and constants are proved from these definitions.



The corrected approximation includes the whole free-OU response. Its exact defect equations lead to cubic velocity and quartic position errors, which combine into the second-order transformed local-error bound.


\section{\texorpdfstring{Build damped integrals and continuous paths}{Build damped integrals and continuous paths}}


Define \(D_\gamma F(t)=\int_0^te^{-\gamma(t-s)}F(s)\,ds\). It is continuous for a continuous \(F\): factor \(\exp(-\gamma t)\) outside the primitive of \(\exp(\gamma s)F(s)\), then apply the fundamental theorem of calculus to that primitive. The general primitive continuity lemma also supports all nested integrals used below.



The damped integrals are defined as Bochner integrals of continuous Hilbert-valued paths. Their continuity and differentiation identities follow from the fundamental theorem and the exponential kernel. These properties ensure the subsequent defect equations concern functions.


\section{\texorpdfstring{The corrected approximation and its two defects}{The corrected approximation and its two defects}}


The corrected velocity is \(U(t)-\psi_1(t)G(0)-H_0\int_0^t e^{-\gamma(t-s)}\int_0^s U(r)\,dr\,ds\). The Taylor defect integrates \(G(s)-G(0)-H_{0}(X(s)-X(0))\) against the damping kernel. The memory defect applies \(H_{0}\) to the damped integral of the primitive of the damped force integral. These definitions retain the entire free-OU path, including its stochastic Hessian correction.



The corrected approximation includes the frozen force, its frozen Hessian response, and all corresponding stochastic terms. Subtracting it from the exact path defines velocity and position defects with zero initial values. Those initial values are important for their higher-order time bounds.


\section{\texorpdfstring{The local-error identity from the exact equations}{The local-error identity from the exact equations}}


From \(V=U-K_\gamma G\) and \(X(s)-X(0)=\int_0^sV(r)\,dr\), obtain \(\int_0^s U(r)\,dr=X(s)-X(0)+\int_0^s(K_\gamma G)(r)\,dr\). Substitute this identity into the correction, commute the bounded linear Hessian through the Bochner integral, and separate the frozen constant force using \(\int \exp(-\gamma (t-s))ds=\psi _{1}(t)\). Algebra then proves \(\widehat V-V\) equals the Taylor defect minus the memory defect. The result assumes exact mild and kinematic equations, rather than an error inequality.



Insert the mild and kinematic equations and cancel the shared responses. The remaining integrand separates the quadratic Taylor remainder from the deviation controlled by the frozen linear dynamics. This identity separates the two terms to be estimated.


\section{\texorpdfstring{Integration of the polynomial bounds with exact constants}{Integration of the polynomial bounds with exact constants}}


For forward time and \(\gamma \ge 0\), the damping kernel lies in (0,1]. If \(\|F(s)\|\le Cs^{n}\), the norm of its damped integral is at most \(\int _{0}^{t}Cs^{n}ds=Ct^{n+1}/(n+1)\). The undamped case follows at \(\gamma =0\). Applying this rule three times to a score bound \(G_{0}\) gives the memory norm at most \(\|H_{0}\|G_{0}t^{3}/6\).



Integrate the polynomial upper bounds exactly rather than replacing every time by the endpoint. Integrals such as \(\int_0^h s^2ds\) and their additional primitives produce the cubic and quartic denominators. These constants feed the final transformed local-error bound.


\section{\texorpdfstring{Obtain cubic velocity and quartic position errors}{Obtain cubic velocity and quartic position errors}}


A Taylor forcing bound \(Rs^{2}\) integrates to \(Rt^{3}/3\). Adding the memory estimate yields velocity error at most \((2R+\|H_{0}\|G_{0})t^{3}/6\). The corrected position is \(X(0)\) plus the integral of the corrected velocity. Its difference from the exact position is therefore the integral of the velocity error, giving \((2R+\|H_{0}\|G_{0})t^{4}/24\). Continuity of the corrected velocity justifies this subtraction of integrals.



A quadratic force remainder contributes a cubic velocity error after one integration. Position is an additional integral and hence has a quartic error. The bounds also account for the linear feedback generated by the frozen Hessian.


\section{\texorpdfstring{The paper’s quarter-constant in transformed coordinates}{The paper’s quarter-constant in transformed coordinates}}


Write \(B=Ct^{3}\). The component estimates give \(\|v\|\le B/6\) and \(\|x\|\le Bt/24\). Under \(\gamma t\le 1/5\), the first transformed coordinate satisfies \(\|v+\gamma x\|\le 7B/40\); the second has norm \(\|v\|\). Squaring and adding gives an energy below \((B/4)^{2}\), hence transformed norm at most \(B/4\). Applying this exact component lemma with \(C=2R+\|H_{0}\|G_{0}\) proves the local phase estimate. For stationary kinetic velocities and score moments, \(C\) becomes \(M_{2}\sqrt{p^{2}+2p}+M\sqrt{Mp}\).



The transformed norm weights velocity and position differently. The step restriction lets the higher-order position contribution fit within the prescribed cubic bound. This is the local second-order estimate that will become a quadratic global bias after geometric accumulation.


\end{document}
